% GENERATED FILE. Edit canonical lessons, then run npm run generate:books.
% canonical-source-hash: fnv1a64:7de4896d19ea2914
% canonical-lessons: RU-C193-shlyapa, RU-C193-remen, RU-C193-karman, RU-C193-pugovitsa, RU-C193-sergi

\chapter{Brimmed hat, Belt, Pocket, Button, Earrings}
\label{ch:ru-brimmed-hat-belt-pocket-button-earrings}
\hlchaptermodality{193}

\begin{chapteropening}
\textbf{By the end of this chapter:} \emph{I can say a brimmed hat, a belt, a pocket, a button and earrings.}

Complete the last lesson of chapter 193: I can say a brimmed hat, a belt, a pocket, a button and earrings.
\end{chapteropening}

\section[\texorpdfstring{\ru{шляпа} (shlyápa)}{шляпа (shlyápa)}]{\ru{шляпа} (shlyápa) — a brimmed hat}

\label{lesson:RU-C193-shlyapa}



\begin{quote}
Before the new one: say the Russian for a ruler, then the Russian for a folder.
\end{quote}

\subsection*{You'll want to know: \ru{шляпа}}
\textbf{\ru{шляпа}} (\emph{shlyápa}) — \textquotedblleft{}a brimmed hat\textquotedblright{}.

\begin{cousinweb}
\textbf{\ru{шляпа}} — recorded from 1589; it comes from Bavarian \emph{Schlappe}, \textquotedblleft{}a cap\textquotedblright{}, Middle High German \emph{slappe}, \textquotedblleft{}a soft, drooping hat\textquotedblright{} (Vasmer).
\end{cousinweb}

\begin{grammarlens}[title={What you've built}]
One more word for this chapter.
\end{grammarlens}

\subsection*{Your turn}
\begin{itemize}
\raggedright
  \item \emph{Say it:} \emph{shlyápa}
  \item \emph{Say it:} \emph{shlyápa}, once more
  \item \emph{Say it:} say \emph{shlyápa}
  \item \emph{Recall:} say the Russian for a ruler, then the Russian for a folder, then say \emph{shlyápa} again
\end{itemize}

\begin{tcolorbox}[breakable,colback=teal!4,colframe=teal!35!black,title={Before you move on}]
What does \ru{шляпа} mean? (\textquotedblleft{}A brimmed hat\textquotedblright{}.) Say it once more.
\end{tcolorbox}
\section[\texorpdfstring{\ru{ремень} (remén)}{ремень (remén)}]{\ru{ремень} (remén) — a belt}

\label{lesson:RU-C193-remen}



\begin{quote}
Before the new one: say the Russian for a folder, then the Russian for a brimmed hat.
\end{quote}

\subsection*{You'll want to know: \ru{ремень}}
\textbf{\ru{ремень}} (\emph{remén}) — \textquotedblleft{}a belt\textquotedblright{}.

\begin{cousinweb}
\textbf{\ru{ремень}} $\leftarrow$ Proto-Slavic \emph{*remy}, \emph{*remene} — inherited. Vasmer holds that a loan from Old High German \emph{riumo}, \textquotedblleft{}a strap\textquotedblright{}, is impossible on the evidence of Old Church Slavonic and old word-building.
\end{cousinweb}

\begin{grammarlens}[title={What you've built}]
One more word for this chapter.
\end{grammarlens}

\subsection*{Your turn}
\begin{itemize}
\raggedright
  \item \emph{Say it:} \emph{remén}
  \item \emph{Say it:} \emph{remén}, once more
  \item \emph{Say it:} say \emph{remén}
  \item \emph{Recall:} say the Russian for a folder, then the Russian for a brimmed hat, then say \emph{remén} again
\end{itemize}

\begin{tcolorbox}[breakable,colback=teal!4,colframe=teal!35!black,title={Before you move on}]
What does \ru{ремень} mean? (\textquotedblleft{}A belt\textquotedblright{}.) Say it once more.
\end{tcolorbox}
\section[\texorpdfstring{\ru{карман} (karmán)}{карман (karmán)}]{\ru{карман} (karmán) — a pocket}

\label{lesson:RU-C193-karman}



\begin{quote}
Before the new one: say the Russian for a brimmed hat, then the Russian for a belt.
\end{quote}

\subsection*{You'll want to know: \ru{карман}}
\textbf{\ru{карман}} (\emph{karmán}) — \textquotedblleft{}a pocket\textquotedblright{}.

\begin{cousinweb}
\textbf{\ru{карман}} — Old Russian already had the name \emph{Kormanŭ}. Turkic \emph{karman}, \textquotedblleft{}a pocket, a purse\textquotedblright{}, was once taken as the source; Radloff saw it as a loan from Russian, and Räsänen again argued for a Turkic origin (Vasmer).
\end{cousinweb}

\begin{grammarlens}[title={What you've built}]
One more word for this chapter.
\end{grammarlens}

\subsection*{Your turn}
\begin{itemize}
\raggedright
  \item \emph{Say it:} \emph{karmán}
  \item \emph{Say it:} \emph{karmán}, once more
  \item \emph{Say it:} say \emph{karmán}
  \item \emph{Recall:} say the Russian for a brimmed hat, then the Russian for a belt, then say \emph{karmán} again
\end{itemize}

\begin{tcolorbox}[breakable,colback=teal!4,colframe=teal!35!black,title={Before you move on}]
What does \ru{карман} mean? (\textquotedblleft{}A pocket\textquotedblright{}.) Say it once more.
\end{tcolorbox}
\section[\texorpdfstring{\ru{пуговица} (púgovitsa)}{пуговица (púgovitsa)}]{\ru{пуговица} (púgovitsa) — a button}

\label{lesson:RU-C193-pugovitsa}



\begin{quote}
Before the new one: say the Russian for a belt, then the Russian for a pocket.
\end{quote}

\subsection*{You'll want to know: \ru{пуговица}}
\textbf{\ru{пуговица}} (\emph{púgovitsa}) — \textquotedblleft{}a button\textquotedblright{}.

\begin{cousinweb}
\textbf{\ru{пуговица}} $\leftarrow$ Old Russian \emph{pugŭvĭ}, \emph{\ru{пугы}}, \textquotedblleft{}a button\textquotedblright{}, with \emph{-\ru{ица}}. It is usually compared with Latvian \emph{puoga}, \textquotedblleft{}a button\textquotedblright{}, and Sanskrit \emph{puñja-}, \textquotedblleft{}a lump\textquotedblright{} (Vasmer), though Persson separated the Sanskrit.
\end{cousinweb}

\begin{grammarlens}[title={What you've built}]
One more word for this chapter.
\end{grammarlens}

\subsection*{Your turn}
\begin{itemize}
\raggedright
  \item \emph{Say it:} \emph{púgovitsa}
  \item \emph{Say it:} \emph{púgovitsa}, once more
  \item \emph{Say it:} say \emph{púgovitsa}
  \item \emph{Recall:} say the Russian for a belt, then the Russian for a pocket, then say \emph{púgovitsa} again
\end{itemize}

\begin{tcolorbox}[breakable,colback=teal!4,colframe=teal!35!black,title={Before you move on}]
What does \ru{пуговица} mean? (\textquotedblleft{}A button\textquotedblright{}.) Say it once more.
\end{tcolorbox}
\section[\texorpdfstring{\ru{серьги} (sérgi)}{серьги (sérgi)}]{\ru{серьги} (sérgi) — earrings}

\label{lesson:RU-C193-sergi}



\begin{quote}
Before the new one: say the Russian for a pocket, then the Russian for a button.
\end{quote}

\subsection*{You'll want to know: \ru{серьги}}
\textbf{\ru{серьги}} (\emph{sérgi}) — \textquotedblleft{}earrings\textquotedblright{}.

\begin{cousinweb}
\textbf{\ru{серьги}} is the plural of \emph{\ru{серьга}}, \textquotedblleft{}an earring\textquotedblright{}. One view reshapes Old Russian \emph{\ru{усерязь}}, \textquotedblleft{}an earring\textquotedblright{}, but Vasmer finds the \emph{\ru{г}} hard to explain and prefers a loan from an Old Chuvash word for \textquotedblleft{}a ring\textquotedblright{}; Räsänen, in Trubachev's note, derives it from Chagatai and other Turkic \emph{isirga}, \textquotedblleft{}ear pendants\textquotedblright{}.
\end{cousinweb}

\begin{grammarlens}[title={What you've built}]
One more word for this chapter.
\end{grammarlens}

\subsection*{Your turn}
\begin{itemize}
\raggedright
  \item \emph{Say it:} \emph{sérgi}
  \item \emph{Say it:} \emph{sérgi}, once more
  \item \emph{Say it:} say \emph{sérgi}
  \item \emph{Recall:} say the Russian for a pocket, then the Russian for a button, then say \emph{sérgi} again
\end{itemize}

\begin{tcolorbox}[breakable,colback=teal!4,colframe=teal!35!black,title={Before you move on}]
What does \ru{серьги} mean? (\textquotedblleft{}Earrings\textquotedblright{}.) Say it once more.
\end{tcolorbox}
